\chapter*{Abstract}
\addcontentsline{toc}{chapter}{Abstract}

\noindent
Mercado Libre sellers typically detect a trending product only once it is already visible in rankings or social media: by then, supply has already grown and the margin of opportunity has already shrunk. A survey of 120 users confirms the scale of the problem (88.9\,\% report detecting trends late or only sometimes on time), and three in-depth interviews provide qualitative depth, including a real case of estimated profitability loss in a 3-to-1 ratio for not having had the signal two weeks earlier.

This work presents a system that detects emerging products early on Mercado Libre Argentina by building its own historical dataset, obtained exclusively through the platform's official API. Over that history, the system computes a trend score made up of five weighted components (appearance frequency, persistence, catalog position, price stability across sellers, and saturation by number of active sellers), and trains a supervised machine learning model (Random Forest) that estimates the probability that a given trend will hold, reaching a recall of 0.807 and an AUC-ROC of 0.903 on the available history.

The system runs autonomously in production, with data capture automated several times a day, free and paid user account tiers, and real email alerts for followed products. As of this report, it tracks more than 270 products with over 70 days of continuous history. The work also includes a business model with quantified economic-financial viability and an independent brand identity (Tendar), decoupled from the platform it operates on. The results confirm that it is possible to anticipate a product's growth before the trend becomes evident through conventional channels, by combining auditable descriptive metrics with a predictive model trained on real data.

\vspace{0.5cm}
\noindent\textbf{Keywords:} trend detection, e-commerce, supervised machine learning, time series, Mercado Libre.